\documentclass[10pt,a4paper]{amsart}
\usepackage{amsmath,amssymb,amsthm,graphicx}
\usepackage{verbatim}
\newenvironment{reference}{\comment}{\endcomment}
%\newenvironment{reference}{}{}
\newenvironment{slogan}{\comment}{\endcomment}
\newenvironment{history}{\comment}{\endcomment}

% For commutative diagrams we use Xy-pic
\usepackage[all]{xy}

% We use 2cell for 2-commutative diagrams.
\xyoption{2cell}
% Theorem environments.
%
\theoremstyle{plain}
\newtheorem{theorem}[subsection]{Theorem}
\newtheorem{proposition}[subsection]{Proposition}
\newtheorem{lemma}[subsection]{Lemma}

\theoremstyle{definition}
\newtheorem{definition}[subsection]{Definition}
\newtheorem{example}[subsection]{Example}
\newtheorem{exercise}[subsection]{Exercise}
\newtheorem{situation}[subsection]{Situation}

\theoremstyle{remark}
\newtheorem{remark}[subsection]{Remark}
\newtheorem{remarks}[subsection]{Remarks}

\numberwithin{equation}{subsection}

% Macros
%
\def\lim{\mathop{\mathrm{lim}}\nolimits}
\def\colim{\mathop{\mathrm{colim}}\nolimits}
\def\Spec{\mathop{\mathrm{Spec}}}
\def\Hom{\mathop{\mathrm{Hom}}\nolimits}
\def\Ext{\mathop{\mathrm{Ext}}\nolimits}
\def\SheafHom{\mathop{\mathcal{H}\!\mathit{om}}\nolimits}
\def\SheafExt{\mathop{\mathcal{E}\!\mathit{xt}}\nolimits}
\def\Sch{\mathit{Sch}}
\def\Mor{\mathop{\mathrm{Mor}}\nolimits}
\def\Ob{\mathop{\mathrm{Ob}}\nolimits}
\def\Sh{\mathop{\mathit{Sh}}\nolimits}
\def\NL{\mathop{N\!L}\nolimits}
\def\CH{\mathop{\mathrm{CH}}\nolimits}
\def\proetale{{pro\text{-}\acute{e}tale}}
\def\etale{{\acute{e}tale}}
\def\QCoh{\mathit{QCoh}}
\def\Ker{\mathop{\mathrm{Ker}}}
\def\Im{\mathop{\mathrm{Im}}}
\def\Coker{\mathop{\mathrm{Coker}}}
\def\Coim{\mathop{\mathrm{Coim}}}

% Boxtimes
%
\DeclareMathSymbol{\boxtimes}{\mathbin}{AMSa}{"02}

%
% Macros for moduli stacks/spaces
%
\def\QCohstack{\mathcal{QC}\!\mathit{oh}}
\def\Cohstack{\mathcal{C}\!\mathit{oh}}
\def\Spacesstack{\mathcal{S}\!\mathit{paces}}
\def\Quotfunctor{\mathrm{Quot}}
\def\Hilbfunctor{\mathrm{Hilb}}
\def\Curvesstack{\mathcal{C}\!\mathit{urves}}
\def\Polarizedstack{\mathcal{P}\!\mathit{olarized}}
\def\Complexesstack{\mathcal{C}\!\mathit{omplexes}}
% \Pic is the operator that assigns to X its picard group, usage \Pic(X)
% \Picardstack_{X/B} denotes the Picard stack of X over B
% \Picardfunctor_{X/B} denotes the Picard functor of X over B
\def\Pic{\mathop{\mathrm{Pic}}\nolimits}
\def\Picardstack{\mathcal{P}\!\mathit{ic}}
\def\Picardfunctor{\mathrm{Pic}}
\def\Deformationcategory{\mathcal{D}\!\mathit{ef}}
\usepackage{cleveref}
\crefname{originalreference}{original reference}{original references}
\Crefname{originalreference}{Original reference}{Original references}
\makeatletter
\newlabel{lemma-proper-base-change-in-terms-of-injectives}{{0A4B}{source}{}{sourcepointer-1}{}}
\newlabel{lemma-proper-base-change-in-terms-of-injectives@cref}{{[originalreference][0][]0A4B}{[1][0][]source}{}{}{}}
\makeatother
\crefname{originallemma}{original lemma}{original lemmas}
\Crefname{originallemma}{Original lemma}{Original lemmas}
\makeatletter
\newlabel{morphisms-lemma-base-change-finite}{{01WL}{source}{}{sourcepointer-2}{}}
\newlabel{morphisms-lemma-base-change-finite@cref}{{[originallemma][0][]01WL}{[1][0][]source}{}{}{}}
\makeatother
\crefname{originallemma}{original lemma}{original lemmas}
\Crefname{originallemma}{Original lemma}{Original lemmas}
\makeatletter
\newlabel{morphisms-lemma-finite-proper}{{01WN}{source}{}{sourcepointer-3}{}}
\newlabel{morphisms-lemma-finite-proper@cref}{{[originallemma][0][]01WN}{[1][0][]source}{}{}{}}
\makeatother
\crefname{originalreference}{original reference}{original references}
\Crefname{originalreference}{Original reference}{Original references}
\makeatletter
\newlabel{proposition-finite-higher-direct-image-zero}{{03QP}{source}{}{sourcepointer-4}{}}
\newlabel{proposition-finite-higher-direct-image-zero@cref}{{[originalreference][0][]03QP}{[1][0][]source}{}{}{}}
\makeatother
\title{Proper base change holds for finite morphisms}
\author{The Stacks Project Authors}
\date{stacks-project:etale-cohomology.tex:local-snapshot-20261011; GNU FDL 1.2 or later}
\begin{document}
\maketitle
\noindent\textit{Editorial note.} All delivered mathematical text is the author's own, copied from the pinned Stacks Project chapter file etale-cohomology.tex; no statement or proof step was rewritten, reordered or removed.\par\medskip
\noindent\textit{Editorial note.} The counted proof is the author's complete argument: the proof is written out in full: the two geometric inputs from the classification of finite morphisms and the combination of the base change criterion with the vanishing of the higher direct images.\par\medskip
\noindent\textit{Editorial note.} The label lemma-finite names the counted statement, and the proof environment immediately following it in the source is delivered verbatim as its proof.\par\medskip
\noindent\textit{Editorial note.} The counted statement is carried by the author's own lemma environment in the chapter source.\par\medskip
\noindent\textit{Editorial note.} The statement retains the author's slogan environment; the preamble spans therefore include the native definitions of the comment environments (native preamble lines 8-13), and the text is delivered verbatim.\par\medskip
\noindent\textit{Editorial note.} The preamble is otherwise the chapter's own native preamble excerpt already used for this source: the Xy-pic setup (native preamble lines 14-18) and the native theorem declarations with the mathematical macros (native preamble lines 154-217).\par\medskip
\noindent\textit{Allowed standard prerequisites (editorial).}\begin{itemize}
\item Base change reformulated through injective sheaves. Same chapter, the lemma identifying the base change property with the vanishing of the higher direct images of injective torsion sheaves after base change. Supplies the criterion applied at the end of the proof.
\item Finiteness of a finite morphism and of its base changes. Morphisms, the lemmas stating that a finite morphism is proper and that the base change of a finite morphism is finite. Supplies the two geometric inputs of the proof.
\item Vanishing of the higher direct images of a finite morphism. Same chapter, the proposition stating that a finite morphism has vanishing higher direct images. Supplies the vanishing needed by the criterion.
\end{itemize}\medskip
\ifdefined\hypertarget\hypertarget{sourcepointer-1}{}\fi
\noindent\textit{Source reference.} Original reference 0A4B: sources/stacks-etale-cohomology/pinned-original-source.tex, source lines 18360--18360. Prerequisite pointer invoked at proof line(s) 18533; it names a standard fact of the same chapter that the counted argument uses, and every step of the counted proof built on it is retained verbatim in the delivered body.\par\medskip
\ifdefined\hypertarget\hypertarget{sourcepointer-2}{}\fi
\noindent\textit{Source reference.} Original lemma 01WL: sources/stacks-chow/tags.txt, source lines 2486--2486. Prerequisite pointer invoked at proof line(s) 18531; it names a standard fact of a sibling chapter that the counted argument uses, and every step of the counted proof built on it is retained verbatim in the delivered body.\par\medskip
\ifdefined\hypertarget\hypertarget{sourcepointer-3}{}\fi
\noindent\textit{Source reference.} Original lemma 01WN: sources/stacks-chow/tags.txt, source lines 2488--2488. Prerequisite pointer invoked at proof line(s) 18529; it names a standard fact of a sibling chapter that the counted argument uses, and every step of the counted proof built on it is retained verbatim in the delivered body.\par\medskip
\ifdefined\hypertarget\hypertarget{sourcepointer-4}{}\fi
\noindent\textit{Source reference.} Original reference 03QP: sources/stacks-etale-cohomology/pinned-original-source.tex, source lines 8245--8245. Prerequisite pointer invoked at proof line(s) 18535; it names a standard fact of the same chapter that the counted argument uses, and every step of the counted proof built on it is retained verbatim in the delivered body.\par\medskip
\begin{lemma}
\label{lemma-finite}
\begin{slogan}
Proper base change for \'etale cohomology holds for finite morphisms.
\end{slogan}
Let $f : X \to Y$ be a finite morphism of schemes.
Then cohomology commutes with base change for $f$.
\end{lemma}

\begin{proof}
Observe that a finite morphism is proper, see
Morphisms, Lemma \ref{morphisms-lemma-finite-proper}.
Moreover, the base change of a finite morphism is finite, see
Morphisms, Lemma \ref{morphisms-lemma-base-change-finite}.
Thus the result follows from
Lemma \ref{lemma-proper-base-change-in-terms-of-injectives}
combined with
Proposition \ref{proposition-finite-higher-direct-image-zero}.
\end{proof}

\end{document}
